% Figure: the tube-vibration margin, plotting the vortex-shedding frequency against
% shell-side crossflow velocity as a rising line, with a tube natural-frequency band
% drawn flat; the resonance danger zone is where the shedding line crosses the band,
% and the fluid-elastic instability velocity is a hard vertical ceiling to the right.
\begin{figure}[htbp]
\centering
\begin{tikzpicture}
\begin{axis}[
  width=12cm, height=7.5cm,
  xlabel={shell-side crossflow velocity $u$ (m\,s$^{-1}$)},
  ylabel={frequency (Hz)},
  xmin=0, xmax=5, ymin=0, ymax=60,
  axis lines=left,
  legend pos=north west,
  legend cell align=left,
  tick label style={font=\scriptsize},
  label style={font=\small},
  legend style={font=\scriptsize},
  samples=60,
  domain=0:5,
  clip=false,
]
% shedding frequency f_s = St u / d_o, with St=0.2, d_o=0.019 m -> slope ~10.5 Hz per m/s
\addplot[engred, very thick] {10.5*x};
\addlegendentry{vortex-shedding $f_s = \mathrm{St}\,u/d_o$}
% tube natural-frequency band, drawn as a shaded horizontal strip 28 to 34 Hz
\addplot[englime!25, fill, draw=none] coordinates {(0,28) (5,28) (5,34) (0,34)} \closedcycle;
\addlegendentry{tube natural-frequency band}
\addplot[englime, dashed, thick, domain=0:5] {28};
\addplot[englime, dashed, thick, domain=0:5] {34};
% resonance crossing markers (f_s hits 28 and 34 Hz)
\addplot[only marks, mark=*, engblue] coordinates {(2.67,28) (3.24,34)};
\node[engblue, font=\scriptsize, anchor=south east] at (axis cs:2.6,29) {resonance onset};
% fluid-elastic instability ceiling as a vertical line at u = 4.2 m/s
\addplot[engamber, very thick] coordinates {(4.2,0) (4.2,60)};
\node[engamber, font=\scriptsize, rotate=90, anchor=south] at (axis cs:4.05,10) {fluid-elastic instability};
% design point well below both hazards
\addplot[only marks, mark=square*, enggray] coordinates {(1.5,15.75)};
\node[enggray, font=\scriptsize, anchor=north west] at (axis cs:1.55,15) {design point};
\end{axis}
\end{tikzpicture}
\caption[Tube-vibration margin against crossflow velocity]{The vibration margin of a
tube bundle read against the shell-side crossflow velocity. The vortex-shedding
frequency rises linearly with velocity, $f_s = \mathrm{St}\,u/d_o$, drawn as the
sloping line, while each tube carries a band of natural frequencies set by its
unsupported span between baffles, drawn as the flat shaded strip. Where the shedding
line enters the natural-frequency band the tubes resonate and grow amplitude until
they fatigue at the baffle holes or clash with neighbours, the region marked
resonance onset. A second and harder ceiling is the fluid-elastic instability
velocity, drawn as the vertical line, above which the whole array goes unstable at
once regardless of shedding resonance. A sound design sits at the marked point, well
below the natural-frequency band and far below the instability ceiling, and the
lever that moves both hazards is the baffle spacing, which sets the unsupported span
and thus the natural frequency and simultaneously fixes the crossflow velocity.}
\label{fig:vol04ch06:vibration-margin}
\end{figure}
